\quellseite{519}
\chapter[§ 35. THE METHOD IN THE CATECHISM]{§ 35. THE METHOD OF INSTRUCTION IN THE CATECHISM}

Compare the literature in ch. 14, 26, and 32. *G. v. Zezschwitz, ii 2: \emph{Die Katechese od. die kirchliche Unterrichtsmethode}, 1859. *F. Lange, pp. 167–241. *A. Eckert, pp. 115–163. J. Gottschick, pp. 179–186. J. Berndt, pp. 99–105. R. Kabisch, pp. 246–247. J. Steinbeck, pp. 201–218. M. Reu, \emph{Quellen z. Gesch. d. kirchl. Unterrichts}, 1904 ff. O. Schoenhuth, \emph{Methodenlehre}, 1903. *B. B. Hahn, \emph{Luther's Katechismus erklärt aus bibl. Theologie}, 1891. A. Scheller, \emph{Die pädagogische Bedeutung d. luth. I. Hauptstueckes}, 1905. F. Niebergall, \emph{Die paulinische Erlösungslehre i. Konfirmandenunterricht}, 1908. K. Eger, \emph{Taufe und Abendmahl im kirchl. Unterricht der Gegenwart}, 1911. H. Matthes, \emph{Der 2. Artikel auf bibl. geschichtl. Grundlage}, 1913. D. Vorwerk, \emph{Gebet und Gebetserziehung}, 1913. PRACTICAL HELPS: J. Nissen, \emph{Unterredungen ueber den kl. Kat. Luthers}, 1885. F. W. Schuetze, \emph{Entwürfe und Katechesen zu Luthers kl. Katechismus}, 3 parts, 1888. K. Ender, \emph{Handbuch zum kl. Katechismus L's fuer Lehrer und Pfarrer}, 1874. L. Wangemann, \emph{Einführung i. d. Verständnis d. Luth. Katech.}, 3 parts, 1880 ff. H. Cremer, \emph{Unterweisung i. Christentum nach L's kl. Katech.}, 1883. O. Zuck, \emph{Katechesen ueber L's Katech.}, 1883. G. v. Zezschwitz, \emph{Christenlehre}, ²1883 ff. J. Chr. G. Schumann, \emph{Handbuch d. Katechismusunterrichts}, 3 parts, 1894 ff. F. H. Kahle, \emph{Der kl. Kat. L's, anschaulich, kurz u. einfach erklärt}, 1896. K. Boehruecker, \emph{Der Katechismusunterricht}, 1895. A. Nebe, \emph{Der kl. Kat. ausgelegt aus Luthers katech. Werken}, 1891. Fr. Oehme, \emph{Die 5 Hauptstuecke d. luth. Katech. katechetisch bearb.}, 1891. Ziethe, Rohden, Heydt, \emph{Die unterrichtl. Behandlg. d. 6. Gebots in der Schule}, ²1894. A. Gruelich, \emph{Skizzen z. unterrichtl. Behandlung d. kl. Katech. L's}, 1899. Th. Hardeland, \emph{Die katech. Behandlung d. kl. Katech. L's in Unterredungen}, 1899. *J. H. A. Fricke, Handbuch d. Katechismusunterrichts, 3 parts, ²1899. *S. Bang, \emph{Katechetische Bausteine zu christozentrischer Behandlung des 1. Hauptstueckes}, ²1901. G. Heimerdinger, \emph{Praeparationen f. d. Katechismusunterricht}, 1901. W. Beyschlag, \emph{Christenlehre auf Grund d. kl. Kat.}, ²1903. *R. Stange, Der Katechismusunterricht, 3 parts, ²1903–08. M. v. Nathusius, \emph{Die christliche Lehre nach L's kl. Kat.}, 1904. G. Kaelker, \emph{Der Katechismusstoff f. d. 5–8. Schuljahr}, 1904. A. Eckert, \emph{Der Katechismusstoff d. Konfirmandenunterrichts in synthetischem Gang}
\quellseite{520}
\emph{u. psychologischer Stoffordnung}, 1905. A. u. F. Falcke, \emph{Praeparationen z. Unterricht i. L’s kl. Katech.}, ⁵1906. *R. Steinmetz, \emph{Katechismusgedanken; Beitrag z. katechet. Behandlung d. 5 Hauptstuecke i. Kirche u. Schule}, 3 parts, 1906–13. Hahne, \emph{Praeparationen f. d. Katechismusunterricht}, ²1909. *A. Eckert, \emph{Kinderkatechismus f. Schule u. Kirche, I. u. 2. Hauptstueck}, 1909. *J. Kolbe, \emph{Der kl. Katech. L’s in ausgefuehrten Katechesen}, ⁷1911. *K. Eger, \emph{Katechismustafeln}, 1911. *O. Hardeland, ⁵² \emph{Konfirmandenstunden}, ⁵1914. F. Heyn, \emph{Katechismusunterricht in Reukauf u. Heyn, Ev. Religionsunterricht}, ²1913. *K. Eger, \emph{Evangelische \emph{Jugendlehre}; Hilfsbuch z. religioesen Jugendunterweisung nach Luthers Katech.}, ²1912. *Th. Kaftan, \emph{Auslegung d. luth. Kat.}, ⁶1913. *H. J. Schuh, \emph{Catechizations on Luther’s Small Catechism}, 1914. Geo. W. Lose, \emph{Catechism Bible Narratives. A Series of Bible Narratives on the Five Chief Parts of Luther’s Small Catechism}, 1915. ON APPLICATION: *K. Caspari, \emph{Geistliches und Weltliches}, ²³1915. P. v. Zychlinski, \emph{Illustrirende Aussprueche, Sentenzen und Geschichten z. Gottes Wort}, 1900. L. Pestalozzi, \emph{Die christl. Lehre i. Beispielen}, ³1901. Baum, \emph{Erzaehlungen und Beispiele z. kl. Katech. L’s}, ²1907. *J. Besch, \emph{Aus der Lernstube des Lebens}, ²1913. Leipziger Lehrerververein, \emph{Im Strom des Lebens. Narratives on the Catechism} (Columbus, Ohio). *G. Warneck, \emph{Die Mission i. d. Schule}, ¹⁴1912. Th. Schaefer, \emph{Die Innere Mission i. d. Schule}, ⁷1912. O. Koenig, \emph{Die Mission i. Katechismusunterricht}, 1913. F. W. Foerster, Jugendlehre, ⁶⁵1915. Compare also \emph{The King’s Highway Series} by Sneath, Hodges, and Tweedy, and \emph{The Golden Rule Series} by Sneath, Hodges, and Stevens.

True as it is that instruction in Biblical History, if it be of the right sort, is at the same time an introduction into the truths of the Catechism, separate instruction in the latter remains necessary notwithstanding. There exists no other way of giving the young a concise and comprehensive introduction to the faith as it lives and manifests itself in the mature congregation (pp. 310 ff.). The method of such separate instruction in the Catechism is the next subject claiming our attention.

If the aim of religious instruction as a whole has been rightly described on page 280, and if instruction in the Catechism is a part of religious instruction in general, it must contribute its part to the attainment of that aim. Every
\quellseite{521}
method, therefore, is erroneous by which those vital sacred truths fail to find that firm and deep anchorage which is needed for the stimulation of the emotions and of the will. Now, the method in vogue in \emph{Reformation times} does not measure up to the standards which we have set. It is true that all the elements necessary to an efficient method were on hand. Luther had inserted biblical pictures in his Catechism, and most catechists followed his example in their editions of Luther’s Catechism. In the preface Luther had advised teachers to adduce many illustrations from Scripture (p. 114); also in this regard there is no lack of followers, among whom Loener, 1544, and the authors of the Joachimsthal Catechism, 1574, are prominent. The practise prevailed of connecting the Catechism with the Bible and of adducing proof passages from the Bible; for instance, Loener, 1544, Huberinus, 1544, Jos. Opitz, 1583. The hymn was linked to instruction in the Catechism, not only in so far as, in the catechetical period on Sunday, hymns were sung that comported with whatever part of the Catechism happened to be under discussion (a very common practice), but also by adducing illuminating hymn stanzas in connection with the several parts of the Catechism (e.g., G. Walther, 1581). Victorius, in 1591, even published a catechism in which the individual sections were accompanied by a biblical picture, the text of a pertinent Bible story, printed in full, and a number of Bible passages.

This material, however, was not turned to account for the development of an efficient method. It was deemed sufficient that the father should recite the several parts of the Catechism before the family, and the teacher before his class, whereupon they were repeated until retained by the memory. When catechisms with explanations appeared, the explanations were taught in the same way. Camerarius and Lossius improved the method slightly; but they found few
\quellseite{522}
followers. The former would explain the Catechism to the upper classes in coherent discourses in which illustrative stories (taken chiefly from ancient history or invented by himself) were interwoven with the catechetical material, and then address a few summarizing questions to the pupils (cf. Reu, Quellen, i 1, 2 p. 98); the latter often introduced Objectiones, i.e., misgivings to be cleared away by instruction. We must indeed admire the exemplary zeal with which the Catechism was “pounded” into the young and common people, and the unlimited faith in the reconstructive and regenerative power of the Word of God; we must also recognize that some of the teachers succeeded in keeping their catechetical explanations free from specifically dogmatical elements (although the number of such teachers was steadily diminishing). The manner, however, in which the Catechism was offered to children was anything but exemplary. If, notwithstanding, the German people experienced a new birth and if a generation grew up which could endure the terrors of the Thirty Years’ War, such results, instead of demonstrating the correctness of the catechetical method in vogue, should rather be considered as evidence that God may bless the efforts of men who are faithful even though they may use imperfect methods. Moreover, the pronouncing and repeating method was much more expedient in an age when church membership was a matter of course and the consciousness of church authority lived as a power in the hearts of the people to a degree impossible today.

The Magdeburg pastor and educator Andrew Cramer, living toward the end of the sixteenth and beginning of the seventeenth century, is possibly no unbiased witness, but he is probably right in regard to the main point when, in his \emph{Anleitung, wie die Jugend in Gottesfurcht, Kuensten und Sprachen zu unterweisen}, he says concerning religious instruction as practiced in his time: “The Catechism is indeed in use in the schools, but it is taught in a very inefficient manner. A lesson is assigned to the boys, which they have to memorize and cram into their heads. While pronouncing the Catechism to the boys and their repeating it is good enough as far as it goes, it does not go far
\quellseite{523}
enough. No explanation is offered in the form of sermons adapted to children or by means of Bible passages and illustrations; and they are not given an adequate spiritual understanding of the Catechism. Nor is any application to life and Christian practise in evidence. When the boys are about to enter the adult stage, the Catechism is treated in the same manner by means of a foreign tongue: with glosses in a foreign tongue the explanation is made; but all to no purpose, ‘\emph{memoria coacta, sine intellectu, sine praxi}’, Christianity, therefore, has not been planted in the right fashion!”—An element of progress is found in the method of \emph{John Valentine Andreae} (pp. 125, 127). He had his pupils first memorize the text of the Catechism, or of his \emph{Evangelische Kinderlehre aus heiliger, goettlicher Schrift} (1621). That done, however, he would enable them, by a freer treatment, frequent reviews, clear generalizations, sound explanations and skillful application, to reproduce intelligently the words memorized and “to give a ready and reasonable account of their faith”. Cp. his \emph{Theophilus} (edited by V. Fr. Oehler, 1878, pp. 22–25 and pp. 28–49). In the same work we read (p. 128): “A good teacher leads, while a bad one drags; the one sheds light, the other darkness; the one teaches, the other confuses; the one guides, the other pushes; the one stimulates, the other depresses; the one scatters cheer, the other fear; the one builds up, the other tears down. In short, unless the teacher himself be a book, yea, a library, an itinerant museum; unless he himself be a personification of the right method of handling and superintending the task; unless he himself be the embodiment of the genius and method of all language and science, and, in addition, an ornament and a flower of State and Church, he, once for all, is not according to my taste. For to take up books and finish them one after another, to drive and to goad to exertion, to enforce injunctions, rules, and fiats and to hammer them in—that is something anybody can do; but to lay down the main contents of a lesson, to prepare the way to an understanding of it, to make the application, to teach the correct use, to go ahead with a good example, and finally to bring everything into harmony with Christ:—that is the teacher’s true function, that is a task worthy of a Christian, for which all treasures on earth are not sufficient remuneration.” Also the Danzig school superintendent and pastor John Maukisch made an effort to improve the method. In Gemeine Nachricht, wie man die Jugend zu Hause und in den Schulen den Katechismus Lutheri abfragen koenne (1653) he advocated the analyzing method, which attained to dominating influence since the days of Spener.

\quellseite{524}
Neither is the analyzing method (pp. 135, 137, 140) recommended and adopted by Spener and Francke, and later influenced by the philosophy of Wolff, a practicable way to the true goal of education and training. Whether it was employed merely to bring into the light the structure of Luther’s explanation (as was done occasionally even in the sixteenth century), or whether by the use of the question and answer method, the several sentences were made easier of interpretation through grammatical analysis into subject, predicate, object, etc., a remedy may indeed have been found against inattention, but the method remained inadequate. It certainly contains an element of truth, but it does not even succeed in thoroughly stirring the intellectual life of the soul, let alone the feelings and the will, in spite of the fact that, in other respects, Spener and Francke were very much intent upon reaching the emotions and upon the demonstration of religious knowledge through deeds.

\emph{Spener} says in his \emph{Bedenken}: “Examinations should consist not merely of questions and the repetition of the answers found in the Catechism, but in the analysis of each question into many other questions, which, in turn, should be so formulated as to possess explanatory force, so that the young are trained to reflect upon the subject and to answer from their own mind” (vol. i, 631 f.). In the preface to his catechism (1677) he says: “As to the use of this little work, the idea is not to expect a few people merely to commit the questions and answers to memory—a torture of the memory which I would rather warn against than advise. It is my opinion that writings such as this should be used as an aid to the intellect rather than as a burden to the memory. For this reason I require that, by the use of these questions and answers, the subject should be so thoroughly imparted that people can answer from their own mind in their own words; this is better than they should commit to memory the very finest formulas. . . It may be laid down as a general principle among us that the youth entrusted to us should not only understand the truths of the faith with which they have become familiar, but also appreciate them as a divinely established rule of life and as an incentive to true godliness. Occasional exhortations between examinations for the purpose
\quellseite{525}
of adapting the truths learned to life will be found expedient. It would also be exceedingly profitable for the preacher who is occupied with this humble work to train the young to go to Scripture itself for proof by turning to the Bible, or at least the New Testament—a book which they should always have with them. Instruction ought to be given how to understand a passage according to text and context; how to analyze it; how to determine its meaning, and how to use it in proving the truth. This is a fine intellectual exercise and discipline; it results not only in such knowledge of the Bible that any passage can be readily found, but also in attention when the Bible is read and in the earnest desire to miss nothing. The ultimate gain will be that the people’s faith, which, after all, is based not on the Catechism, but on Scripture, becomes well grounded, and that the conviction is wrought in them that their Catechism is really founded upon the Word of God.” — \emph{A. H. Francke}, in his \emph{School Statutes} of 1702, says: “The method of catechization consists 1. in the recitation; 2. in the explanation; 3. in the application. 1. The teacher lets the children recite the part which he intends to treat; 2. he shows them the true meaning of every word of the Catechism, lest they rattle off the catechetical text without making the progress to be expected; 3. he points out to them how the things taught and explained should become for them (a) a solid foundation for their faith and (b) a means of testing and amending their life. All this is not to be imparted by long-winded talks but by the simple method of questions and answers, and with great love, meekness, and kindness. However, the catechetical method of some is beset by faults, which should be noted: 1. there are those who formulate the question in such a way as to require an unvarying yes or no; this has the effect of discouraging the children or so habituating them to shouting yes and no that they lose attention; 2. others ask very few questions when they catechize, but they lecture incessantly and preach regular sermons, which has a tendency to make the child restless; 3. then, it is found that some catechists do not abide by the subject; they pass from one topic to the other and go roving everywhere; but the subject of the catechization is but lightly touched or altogether ignored, whereby it is obscured rather than explained: in these circumstances the children learn nothing whatever with thoroughness; 4. the attempt is also made to arouse the flagging attention of the children by the infliction of bodily punishment; this does more harm than good. . . . Furthermore, a Bible passage should be made clear through questions; for instance, Christ gave Himself for us. Who gave Himself
\quellseite{526}
for us?—For whom did He give Himself?—What did He do for us?—What did He give for us?—When the children have learned to understand a passage in this way, the catechist should induce them to apply it; for instance, Who is it that gave Himself for us? Christ.—For whom did He give Himself? For us. (For me.)—Should we not love a Savior who gave Himself for us? Oh, yes, etc.—This method of imparting knowledge to the children through questions and answers is earnestly recommended to the teachers, not only because this is the best method of explaining a subject, but also because it tends to quiet the unruly minds of the children at a time when there is danger of feeling bored”.

\emph{Loeseke}, in his \emph{Zergliederte Katechismus} (²1758), treats the whole Catechism according to the following plan: “Which article is this? The First Article. What does the First Article treat of? Of creation, Gen. 1:1. What does it say? I believe . . . earth. Who believes? I. What is that you do? I believe, Rom. 1:17. In whom do you believe? In God, John 1:9. What do you call this God? The Father, Eph. 4:6. What kind of maker do you call this God? Almighty. As what is He almighty? As maker. Of what is God the Maker? Of heaven, Ps. 33:6. Of what else? Of earth, Jer. 31:15. How does Luther’s explanation begin? What is meant by this? What do you do when you hear all that is said in the First Article? I believe, 2 Tim. 1:12. What do you believe in regard to your own creation particularly? That God has made me, Eph. 4:24. (Who has made you? Whom has God made? What has God done for you?) Together with what has God made you? Together with all creatures. Together with how many creatures? With all. To whom has God given many blessings and goods through creation? To me. What, for instance, has God given you at your creation? My body. What else? The soul. Gen. 2:7, etc.—Rambach’s manner, indeed, is not quite so dull; a number of valuable hints as to method are given in his booklet, \emph{Der wohlunterrichtete Katechet} (1722); but he does not pass beyond the restrictions of the analytical method. Only in his posthumous work, \emph{Wohlunterwiesener Informator} (1737), noteworthy beginnings of the developing method are found, especially on pp. 140 ff. Here he clearly expresses it as his opinion that Luther’s Catechism cannot be successfully treated with the children until there has been imparted to them 1. some simple preliminary instruction concerning God and divine things, 2. a brief outline of the plan of salvation, and 3. an outline of Biblical History (p. 140–147). Moreover, he clearly recognized that the will of the children, no less than their intellect,
\quellseite{527}
needed correction. “Most informatores”, he says, “who themselves lack information, make a mistake in this matter, and fancy that they have done their duty, and more than their duty when they have imparted to the child the Catechism and a few passages, and when they have taught him to read and to write; but the state of the child’s will appears to concern them precious little. It is all the same to them when they grow up savage and uncouth, their life a field for Satan’s work. In that manner the children are made angels intellectually, but devils according to the state of their will, the aim being not godliness, but knowledge.” What he thereupon says as to the method employed in correcting the will (pp. 162–235), is worthy of attention, but we cannot enter upon it here since he has failed to show how his views concerning catechetical instruction might be instrumental in training the will.

A great step of progress is represented by the Socratic method of the period of rationalism. For the first time now, there appeared a real didactic conversation between the teacher and the pupil, a real development of the unknown from the known (pp. 142 ff.),—a method which really put the intellectual life in motion, at least when employed by its more eminent representatives, especially Dinter. However, the catechesis was too exclusively an intellectual operation, the feelings remaining cold and the will unstirred, quite aside from the substance of instruction, which rarely transcended natural religion. H. Chr. Schwarz: “The insistence upon knowledge was quite in place; but the one-sided way of stressing this principle led to the error that ignorance was worse than wickedness. The heart remaining empty, such catechizing was in the long run no better than the mechanical memorizing previously in vogue.” Cl. Harms: “Question follows question lickety-split; there is discussing, developing, explaining; example and illustration are brought in; punctuation is drawn upon for light, and assumption must do its share, and whatever else is required by catechetics is there, with all of which, if all goes well, logical efficiency is attained and the joy: Now I have caught it. But the adult who already knows
\quellseite{528}
those things and craves a little religious inspiration, cannot stand it. To ask question after question and to call for an answer is not instructing the children. In addition to questions that address themselves to the intellect, such are to be put which, like nails and spears, make for the feelings—the heart. . . Nor should we be content with questions as means to make our way to the heart. Talks, admonitions, entreaties, stories, reading by the teacher, reading by the pupils, prayer and such things are the instruments through which the Kingdom is taken by force.”

Here is an example from the \emph{Unterredungen ueber die fuenf Hauptstuecke des Katechismus} (1806 ff.), by Dinter: Teacher: Has the garden back of your house come into being of itself? Child: No, my grandfather planted it.—Who, then, was there first, the garden or the grandfather who planted it? Grandfather was there first.—Why did that have to be the case? Otherwise he would have been unable to plant it.—What do we call that which is produced, or brought about by someone? The effect.—And what do we call that through the power of which an effect is produced? The cause.—When you think of your grandfather and the garden in this connection, can you use these words? Yes, my grandfather was the cause, and the garden was the effect.—Do you think that the same is the case elsewhere? Yes.—Think of the watch and the watchmaker. The watchmaker was there first.—Think of a similar example! The cabinetmaker was there before the table.—Now think of a case where the effect was there before the cause! That is impossible: if the cause is not there first, it cannot make it at all.—It would be better to say: it cannot bring about an effect: now can you state this in general terms? The cause always comes before the effect.—Now let us apply this to God and the world! God is the cause, and the world is the effect.—What follows from this? God must have been there first, or he could not have made the world.—As soon as I think of a maker of the world, it becomes clear that there must have been a time when the world began to be. Before it was there, there was nothing but God: God made the world: by whom, now, was God made? By nobody: nobody was there but God.—Well, if He was made by nobody, He must have come into being of Himself! That is out of the question; for nothing can come forth by itself.—But is that quite sure? What is meant by saying, You have reason? That I can see the connection between
\quellseite{529}
cause and effect.—Suppose, now, someone would fancy an effect without a cause; what would you call him? Unreasonable.—But coming into being is surely an effect. Therefore there must be a cause for it.—If someone, accordingly, should say: “Something comes into being of itself”, what would you say of him? He speaks unreasonably; where there is an effect, there must be a cause; nothing, therefore, can come into being of itself.—And God? Neither can God have come into being of Himself.—Now let us put all this together: God was there before the world; nobody can have made Him; nor can He have come into being of Himself; but one possibility remains; which might that be? He cannot have come into being at all; the world came into being; that is put also in this way: it had a beginning; as to God, He never came into being.—How can that be best expressed? He has no beginning; he has no end.—Good that you remember this from the lower grade; but there you had to learn many things without knowing the reasons why we should believe that it is so; now your attention is to be called to the reasons, too.

Do you see yonder tree? What do you believe concerning it? Will it be there forever? No, some time it shall be cut down.—Why do you believe that? Because the people need wood, if they are not to freeze.—Now, when someone cuts it down, does the cause of its destruction come from the outside or from itself? It comes from the outside.—Suppose that nobody should cut it down, will it always stand there, and will it never be destroyed? Sometime it would decay and cease to be.—In that case, would the cause of destruction come from without? No, it would come from within.—This is the case with the tree; this is the case with every other thing; when it is destroyed, only one of two things is possible; where may the cause of destruction be found? The ground of destruction is either without or within the thing destroyed.—No third possibility can be imagined; let us apply this to God, and ask the question whether God can ever take an end. If God were to have an end, one of two events would have to take place. Which? The ground of God taking an end would have to be within Him or without Him.—We will speak in connection with God not of destruction, but of extermination. Why? That shall be explained in one of the following conversations; if, now, of two things the one destroys the other, which of the two is bound to be the stronger? The one that does the destroying is bound to be stronger than that which is destroyed.—Now, suppose that God, the living God, God the Creator, should be exterminated by someone, what would have to be the case? The other would have to be stronger than God.—But is
\quellseite{530}
that a possibility? No, nobody is stronger than God; for God is almighty.—What, accordingly, is no living being, no power of nature, able to do? To exterminate God.—The cause of an extermination of God, therefore, cannot come from without; but we assumed a second possibility: could the cause of extermination lie within God as it does within the oak? No, that, too, is impossible.—Why? Suppose that yonder oak tree had a duration of life of five hundred years, could it, in that case, have been there eight hundred years ago? No, it would have decayed a long time ago. How many years ago? Three hundred years ago.—Or could it be possible that it had no beginning at all; that it was there from the beginning? In that case it would have decayed even sooner.—This much, then, must be clear to you: a thing that bears in itself the cause of destruction, cannot have been there from the beginning; why? It would have been destroyed long ago.—Hence, if anything has no beginning, what are you sure of in regard to it? What can it not bear within itself? The ground of destruction.—Apply this to God! God has no beginning; for that reason He cannot bear the ground of extermination within Himself.—There is no ground of extermination within Him or without Him. What do you conclude from that? That no ground whatever of His extermination exists.—Therefore, just as surely as God has no beginning, He has . . . . can you continue? He has no end.—What it is now that you have said of God? He has neither beginning nor end.—Do you know what He is called for that reason? He is called eternal.—Right; sometimes we call Him something else; whatever comes to an end, we call what? Finite.—Him who never takes an end, we call what? Infinite.

Let us once more survey the path which our conversation has taken! What did we first say of cause and effect? The cause is there always first, before the effect; God is the cause of the world; hence God must have been there before the world.—Very well; but that does not say that He was there from eternity; someone may have produced Him. No; for no one was there but Himself.—How did we argue further? He cannot have come into being of Himself; for nothing comes into being of itself.—Now, then, if He has not been produced by anyone else nor come into being of Himself, what must we conclude? He did not come into being at all; He had no beginning.—But whence do we know for certain that He has no end? If anything is to come to an end, the cause must be found within or without it.—Which of the two, now, may be the case with God? Another cannot exterminate God because God is almighty; nor can a ground of extermination be found in God; for He has no beginning.—From all this
\quellseite{531}
follows the divine attribute about which we mean to converse today; which? God is eternal.—That means? He has neither a beginning nor an end.

In the period of a more general \emph{return to the faith of the Reformation} many catechists made the regrettable mistake of resuscitating the old method. In contrast to the Socratic method of the rationalistic period, they would simply impart the catechetical statements to the children by the force of authority, draw them out again by means of questions, demand the proof for them in the form of an occasional Scripture passage, and, possibly tell a story from life in order to shed light upon it or insure its retention by the memory. When they would really employ examples from Biblical History, they did not utilize them fully; least of all did they use them as sources of deduction for the truths of the Catechism. For instance, when a definition of faith had been given and proved as correct by several passages from Scripture, it was added: Thus Abraham, David, the centurion of Capernaum, the king’s officer, the jailor of Philippi, and others, believed. That such proceeding was worthless, is evident. Others recognized the indefeasible element of truth in the Socratic method; they saw that it does not suffice merely to impart and draw forth again by questions; to analyze, and, perhaps afterward, to demonstrate and illustrate—a process during the whole of which the attitude of the children is purely receptive, but that, through the didactic conversation, there must take place a real developing, and thereby a training of the child for self-activity (so Harnisch especially). But which was to be the starting-point for such a developing process? Entertaining the perfectly correct view that instruction in the Catechism has accomplished a great deal when by means of it an understanding of the words of the Catechism has been imparted to the children, some would proceed from the linguistic usage in which the concepts of the Catechism had found expression in the language of the Bible and the vernacular in general (Brieger
\quellseite{532}
particularly). Others would make the Bible passage their starting-point, endeavoring to develop from that the particular Catechism truth to be taught (they were of the opinion that only in that way the “formal principle” of the Reformation could come into its own); Materne, 1853, and, even more, \emph{Crueger}, 1860, represented this method. The latter says plainly: “To adduce Scripture passages as proof material, after the sentence of the Catechism has been discussed, presupposes doubt in the correctness of the doctrine, and for the reason that the passages, for lack of time, are afterward not carefully explained, it confirms such doubt.” As a matter of fact, by that kind of development (especially the type recommended by Crueger, see example below) little was accomplished for the intellectual life of the pupil, and nothing at all for the feelings and the will, let alone for practical life. Everything was abstract rather than concrete and intuitional; the explanation was not gained from the concept sphere of the child. Such a method did not only leave the heart of the child cold; it might even bring about the notion that the Bible is merely a code of doctrines requiring absolute submission. Thereby, however, joy in Scripture is rather choked than fed; and yet, to arouse and nourish it, is one of the very objects of instruction.

Something decidedly better was achieved by G. v. Zezschwitz and all those who, more or less independently, followed the path blazed by him both in theory and practise. At heart in agreement with the truth of Scripture, he traced his steps back to Dinter in regard to form. Repudiating every catechetical method but the developing one, he laid all possible stress upon the elaboration of clear conceptions and precise definitions. He became the father of the positive scientific catechization (Kunstkatechese), which, by a series of severely logical questions, forced the pupil to reach a series of conclusions along a gradual ascent of logical development, and,
\quellseite{533}
by way of summing up these conclusions, the final theme—the exact and precise determination of the particular catechetical truth to be ascertained (cf. Zezschwitz’s Model Catechization on the Ninth and the Tenth Commandments, Katechetik ii, 2, p. 562 ff. and his \emph{Christenlehre}, three vols, ²1883—1888). It cannot be denied that the co-operation of the children was secured in this way to a rare degree, and that their mental powers were subjected to an extraordinary discipline; likewise, that they attained accurate religious knowledge. At the same time it is true that, in point of contents, there was a deplorable overloading by the introduction of material not rooted in the Catechism but dragged in from biblical theology and dogmatics; that an undue amount of attention was given to the formation of conceptions; that through failure to accord all-sided recognition to the intuitional principle only little influence was brought to bear upon the emotional and volitional life; that too much was expected from the average child, and that insufficient scope was given to pupils to give unrestrained expression of their own thoughts.

In his \emph{Versuch einer biblisch-sachlichen und sprachlichen Erklaerung des Kleinen Katechismus} (³1853), \emph{C. F. Brieger} suggests a proceeding something like this: Suppose the term to be explained is “redemption”. The root-word is emo = to purchase. Inasmuch as the purchase of a human being postulates a condition of slavery, redemption means that one who had been a captive or a slave has been purchased for the purpose of restoring him to liberty, the prefix re designating a return to the condition preceding the state of captivity. That which requires redemption is, as it were, in a state of slavery, precluding freedom of movement. With this linguistic conception Holy Scripture agrees, for whenever it speaks of anything from which redemption is required, it is something which cramps and imprisons, such as enemies, distress, sin (Ps. 18:3; 102:20; 130:8). “By the use of such a method,” Brieger thinks, “the conception is built up, as it were, before the very eyes of the pupil, and impressed so much more thoroughly upon his soul because of the part taken by himself in the erection of the structure.”—Crueger, in his Entwurf einer entwickelnden Katechismuslehre (1860; ¹¹1889), explains “with
\quellseite{534}
all cheerfulness and confidence” in the Introduction of the Lord’s Prayer as follows: “According to the words of Luther’s explanation, the dear children ask their dear Father with all cheerfulness and confidence. We are reminded of the passage Rom. 8:15: “Ye received not the spirit of bondage again to fear . . . Abba, father” (8th Trin.). We Christians have received a child-like mind; by this we are prompted to pray: Dear Father! We have not received the mind of a servant, a slave; a slave fears that some evil will be inflicted upon him by his master. When a slave asks anything of his master, he is afraid that some harm may be inflicted upon him on account of his petition. He who fears that harm will befall him on account of his petition, does not pray cheerfully. We pray cheerfully when we are not fearful of harm as a result of our petition. The dear children pray not only with cheerfulness, but also with all confidence (1 John 5:14: “This is the confidence that we have in Him, that, if we ask anything according to His will, He heareth us”). We have the cheerful confidence, the joyful faith, that God will hear our prayer. We pray to God with the cheerful confidence that our prayers will be heard. This cheerful faith is the right kind of confidence when we pray. We pray, accordingly, with all confidence when we pray with the joyful faith that our prayers are heard. Faith in the love of the heavenly Father moves us to pray.”

When \emph{L. Kraussold} proceeded to revise Dinter’s book: \emph{Die vorzüglichsten Regeln der Katechetik} (1801) but instead published his own book on Catechetics from the orthodox standpoint (1843), he became the forerunner of the positive scientific catechization. E.g., the thesis: “The Grace of God in Jesus Christ, is the Christian’s Solace in the Hour of Death” he develops in the following manner: “What does sin merit? Punishment. What punishment has God decreed for sin? Damnation. What is the sensation of the person who looks for punishment? Anxiety. What, therefore, is the sinner bound to feel when he thinks of his sin and punishment? Great anxiety. Whereby could such anxiety be taken from him? (Whence anxiety? When may a man cease from anxiety? When, accordingly, will anxiety have passed away? When will anxiety disappear? Whereby, accordingly, can he be relieved of anxiety?) By being relieved of sin. What does sin merit? The penalty of damnation. What passes away at once when sin is taken away? Punishment. Who punishes? God. Who alone, accordingly, is able to remit the punishment? God. When God remits the penalty of sin, what is that called? To forgive sin. Whereby, accordingly, is anxiety removed?
\quellseite{535}
By the forgiveness of sin. When one is relieved of anxiety, we call that giving comfort. What, therefore, is it that the sinner seeks when he seeks the forgiveness of sin? Comfort. What, in consequence, is the sinner’s comfort? The forgiveness of sin. Has the sinner deserved it? No. What is it called for that reason? Grace. Who has secured for him such grace? Jesus Christ. In whom, therefore, does the sinner find the grace of God? In Christ Jesus. What, therefore, is the Christian’s comfort in his sin? The grace of God in Jesus Christ. When does the Christian stand in need of such grace? Always. Why? Because he is always a sinner. What must he do to secure grace from God? He must ask God. But how long may he ask God for it? As long as he lives. No longer after death? No. What is it that causes him torment if he has not secured grace before death? Damnation. When, therefore, is he nearest to damnation on earth? At death. When is anxiety bound to be the greater, when punishment is near or distant? When it is near. When, therefore, is the sinner’s anxiety bound to be greatest? At death. When, accordingly, does he stand in need of comfort most? In the hour of death. What, therefore, is the Christian’s comfort in the hour of death? The grace of God in Jesus Christ.”

With this model of a scientific catechization before us, we appreciate the views expressed by \emph{Schmarje}. Though himself not averse to catechizing in this manner, he is more moderate than Zezschwitz. While emphasizing the developing question, he repudiates the „\emph{zwingende}“ question of the typical scientific catechization—those severely logical questions whereby the pupil is forced by the teacher to move along the predetermined path of the latter’s reasoning, nor is he a zealous advocate of the “final theme” as the ultimate goal of every catechization. In \emph{Das katechetische-Lehrverfahren auf psychologischer Grundlage}, ²1892, p. 18, he says in regard to the scientific catechization: “The immediate impression of such a conversation is that the pupil can be conceived only as being helplessly drawn toward the goal with invisible tongs. He trots behind the teacher; but he does not know why he follows him, nor whither he follows. The logical consistency of the questions presents no gap for a detour to the right or to the left which might lead astray; but neither does it present an opening through which a glimpse might be obtained of the inner life of the child.” Upon the thesis: “We should fear God”, he himself offers the following model of a catechization conducted by him in Altona with boys ranging from eleven to twelve years and recorded stenographically: “Teacher: You know that Luther began the explanation
\quellseite{536}
of every Commandment, with the sole exception of the First, with the words: We should fear and love God, etc. . . Why does he do that? Pupil: We shall keep the commandments of God when we fear and love Him.—T.: Today we want to explain what it means: We should fear God. Which of the two expressions will it be necessary for us to explain first if we want to understand what is meant by this demand? Pupil: We should know what it means to fear.—T.: Name persons from Biblical History who showed fear. Pupil A.: Moses had fear of God when He appeared to him in the fiery bush. Pupil B.: Paul had fear, when, on the way to Damascus, he suddenly heard the words: Saul, why persecutest thou Me? Pupil C.: Adam and Eve had fear when they had eaten of the forbidden fruit. Pupil D.: The shepherds in the field, when the angels appeared to them. Pupil E.: Mary had fear when she wanted to see the grave of the Lord and saw the angel sitting before it. Pupil F.: The Israelites had fear when Pharaoh pursued them.—Name also some examples from secular history. Pupil A.: Desiderius had fear when he heard that Charlemagne moved upon him with a large army. Pupil B.: The Romans had fear when Hannibal had defeated them.—Can you name also some examples from poetry? . . . Think, e.g., of “The Ring of Polycrates” Amasis feared the envy of the gods.—There we have quite a number of examples of fear. Now tell me: Why did the Israelites have fear when they stood on the Red Sea? They had fear because Pharaoh was behind them; he wanted to take them back.—Why did Adam have fear? He was afraid of punishment.—Why did Jacob have fear when he returned from Laban and was about to meet his brother Esau? He thought that Esau had evil intentions and would take vengeance.—What is it we want to know when we ask “why”? We want to know the ground.—In how many cases may we inquire about the ground of fear? In every case.—Why? There is always some ground for fear.—You named a short time ago a good many examples of fear. Compare the causes of fear with each other! They agree with each other.—Why? When anyone has fear, he is afraid of somebody’s anger.—Is that your opinion? Have you ever had fear? Yes.—Of what, for instance? Of a thunderstorm.—What did you think when the dark clouds massed themselves upon the sky and the lightnings began to leap? I thought that the lightning might strike.—Do you mean to speak of the wrath of a thunderstorm? No.—You see, therefore, that your answer was not correct. Now compare
\quellseite{537}
all the grounds of fear with each other. No answer.—Well, what was Adam afraid of? Of punishment.—Jacob? Of the vengeance of Esau.—Of what is the child afraid during the thunderstorm? Of the lightning-stroke.—Now, if Adam is afraid of punishment, Jacob of the wrath of his brother, and the child of the lightning, you surely must be able to express an opinion when you compare these several grounds of fear. They are different.—That is correct. And yet you contended at first that all the causes of fear were alike. Can you not give the reason for your assertion? Pupil A.: Yes, when one has fear, he thinks of God.—Really? Pupil B.: No, he always thinks of that which he fears.—What does he expect of the persons or things that he is afraid of? He expects that they will do him harm.—What, therefore, is it that one thinks of when he is fearful? The person who has fear always thinks of harm coming to him.—That is very true. But tell me in other words that the grounds for fear, though outwardly different, are alike inwardly. Pupil J.: When anyone has fear, he always thinks of coming harm. Pupil K.: He always thinks of an approaching danger.—What, accordingly, is the reason for fear? The reason for fear is the thought of some unavoidable danger or a threatening disaster.—Good. Express the same thought also in this form: “We have fear when we, etc.”. We have fear when we think of an unavoidable danger or threatening disaster.—Now we know the reason for fear, and we want to remember it (It is written on the blackboard).

“Let us go a step farther. You said some time ago that, on a certain occasion, you had fear of a thunderstorm. How did you know that you had fear. I had that feeling.—You (turning to another pupil), have you, too, ever had fear? How did you know it? I felt it.—What, then, is fear? Fear is a feeling.—What kind of feeling is fear? Fear is a disagreeable feeling.—Who can express the feeling of fear in a different way? Pupil B.: Fear is a feeling of dread. Pupil F.: A feeling that excites you very much.—That is still better. It may happen that, when fear is felt, we become speechless; yes, that we have difficulty in drawing breath. Fear is a depressing, crushing feeling. The answers are all correct. But we want to keep in mind for the present only the last definition. Answer, accordingly, the question fully: What is fear? at the same time telling us how it arises. Fear is a crushing feeling arising from the thought of some unavoidable danger or threatening disaster.
\quellseite{538}
—This sentence, too, we want to remember. (It is repeated a few times and written on the blackboard.)

“Name feelings similar to fear! Pupil: worry and care.—What do you call the feeling that creeps over you when suddenly danger comes upon you? Fright.—What do we feel when the danger coming upon us is not only sudden but at the same time great and overwhelming? Terror.—Summarize what feelings similar to fear we have found? Worry, care, fright, terror, and expectation.—You have named expectation as a feeling similar to fear. When, on Christmas eve, the pine-tree has been decorated by father and mother, while you stand in an adjoining room full of expectation, do you really think that you have fear? No.—Your answer, therefore, was not quite correct. We now know what is meant by fear or what we mean when we say: We fear. What, however, did we designate as the subject for today? We are to learn what is meant when we say: We should fear God.—Give another expression for, We should fear God. We should have fear of God.—Have fear of God! What is God’s relation to our fear? No answer.—You know what I mean when I say: Jacob had fear of Esau; the Israelites had fear of Pharao. What do we mean to say when we speak of fear of God? Pupil A.: God is the one that excites our fear. Pupil B.: God is the ground of our fear.—What had we just recognized as the ground of our fear? The thought of an unavoidable danger or a threatening disaster. And now you contend that God is the ground of your fear? How do you reconcile the two statements? Pupil M.: They cannot be reconciled; the fear of God is something different.—Who is of a different opinion? Pupil C.: They can be reconciled. God, too, can send us disaster.—Explain that more fully! God has the power to do so, and it may be intended as a punishment.—What can be intended as a punishment? The disaster.—What disaster did Adam think of when he hid from God in the garden? He thought of punishment for sin.—In how far can we apply the statement that fear comes from the thought of a threatening danger or disaster, to the fear of God? The fear of God comes from the thought of a threatening danger, because the punishment for our sin is a threatening danger.—What kind of disaster does the godly man consider punishment for sin when he compares it with every other evil? As the greatest evil of all.—Why? Pupil K.: Every other disaster passes; this one lasts forever. Pupil F.: Every other disaster comes upon us in this life, while that one
\quellseite{539}
comes upon us in the other life as well.—Which parable comes to your minds in this connection? The parable of the Rich Man and Poor Lazarus.—Name just one feature of that parable that illustrates the terribleness of the punishment yonder! The rich man said to Abraham: Send Lazarus, that he may dip the tip of his finger in water and cool my tongue; for I am in anguish in this flame.—If the godly man considers punishment for sin the very greatest evil, what will his fear of God be like? Compare it with every other fear! It is greater than any other fear.—Make that clearer! Pupil J.: The godly man has more fear of God than of men. Pupil K.: The godly man fears God above all things.—Give me some examples! When the Christians were persecuted the demand was made upon them to cast aside their faith; but they refused to do it.—But why? They feared God more than any man.—Other examples! Abraham was willing to sacrifice his own son.—What did he prove by that? That he feared God.—Go on! Luther at the diet of Worms, when he was told to recant his writings.—What did he say? I cannot recant. Here I stand; I cannot do otherwise; God help me. Amen!—What do you seek to prove by that example? That Luther feared God more than he did the emperor.—We have become acquainted with one ground of the fear of God. But we must look at this ground somewhat more closely. In what light does man consider God when he fears His punishment? As the judge.—What do you think of a fear of God in which man thinks of God merely as a judge about to inflict punishment? That is a wrong kind of fear of God.—Can you give me another name of such wrong kind of fear of God? Servile fear.—Why? It is the same kind of fear which a servant has of his master.—What is your opinion of the fear of God felt by the Jews? The fear which the Jews had of God was the servile kind.—Why was that? They thought of God merely as a severe judge.—What do you know of the fear which the heathen have of God? The heathen have no fear of God at all.—How is that? The heathen do not believe in God at all, but only in idols.—That's right; in the same sense as in the case of the Jews we cannot speak of the heathen having fear of God. But they fear their gods and idols. What is your opinion in regard to this fear? It is a servile fear.—Prove that. Pupil B.: The heathen performs sacrifices in order to appease the wrath of their gods. Pupil M.: They even sacrifice their own children.—Quite true. Think of the poem: “The Ring of Polycrates.”
\quellseite{540}
There, too, we behold the servile fear of the heathen.—How? The king is to sacrifice his last jewel, lest the gods should do him harm.—Which words of Amasis can we quote as proof of the truth of our statement? The envy of the gods I dread; From birth until he joins the dead, No mortal always lives in glee.—What kind of fear, then, is it which we Christians should feel? A child-like fear.—Assign a reason for it. Pupil C.: God is our father, and we are His children.—What kind of fear this child-like fear is, I want to make plain to you through an example from life. (Here follows an incident taken from life.) The boy had fear. Of what? Of punishment.—That, too. But therewith you have not sufficiently described the ground of his fear. He feared to cause grief to his father.—And why was the thought so painful? He knew that his father loved him.—Now we have brought out two thoughts that we can give as grounds of his fear. Which are they? The boy thinks first of the punishment he has to expect, and, second, of the possibility of grieving his father, who loves him so dearly.—Apply this to the fear of God. But give a complete answer, in this form: The servile fear arises from the thought, etc. The servile fear arises from the thought of the punishment which a person has to expect for his sin; the child-like fear arises from the thought . . .—Well? Remember that we have given a twofold cause of the child-like fear. The child-like fear arises from the thought of the punishment for our sin and of the grief that we caused a loving God by our sin.—Let us remember this statement (it is put down on the blackboard).

“Let us now repeat the main statements that have been made, by finding answers to the following three questions: 1. How does fear arise? 2. What is fear? and 3. How does the true fear of God arise? Fear arises when we think of some unavoidable danger or threatening disaster. Fear is a crushing feeling arising when we think of some unavoidable danger or threatening disaster. The true fear of God arises when we think, on the one hand, of the punishment that we have merited by reason of our sin, and, on the other hand, of God, the loving father, whom we have grieved by our sins.—We have so far considered only the ground of the true fear of God. Let us view next the effect of the true fear of God. What will the servant do who fears the wrath of his master? Pupil G.: He will avoid him. Pupil H.: He will flee from him.—Apply this to the servile fear of God! When a person has a servile fear of God, he will flee from God.—Examples!”
\quellseite{541}
When Cain had slain his brother Abel, he fled to another country. Pupil D.: When Adam had sinned against God he hid from Him.—What effect, now, is the child-like fear likely to have? Think of the boy of whom I told you a short time ago. He went to his father and confessed his shortcoming; he thought of it again and again and avoided sin after that.—What effect, accordingly, has the child-like fear of God upon the disposition and conduct of people? It makes them better.—Quite right. The servile fear of God causes us to flee from God; the child-like fear of God, to flee from sin: we may accordingly say that it sanctifies us. What kind of feeling, accordingly, is the true fear of God? It is a feeling that sanctifies.—Give, therefore, a complete answer to the question: What is the true fear of God?, the ground from which it comes into being to be included in the answer. The true fear of God is a sanctifying feeling arising from the thought that we have sinned, and that we grieve the loving God through our sin.”

This example of Schmarje’s reminds us of von Rohden’s comment upon the scientific catechesis: “As, in the game of croquet, the ball is driven by many blows, both direct and auxiliary, to a small goal through narrow arches, immovably fixed at the outset, so the master of the scientific catechization, by means of many chief questions, and, at every sideward leap of the childish fancy, of even more numerous auxiliary questions, drives the thought of the pupil with unwavering determination through the previously fixed iron gates of his own pre-arranged course of thought to the goal before him—the conception formulated by the book or the teacher, a definition or something like it.” While this criticism is too pungent, it corroborates most of what has been said on the subject.—Zezschwitz, in his model catechesis on the Ninth and the Tenth Commandments, first achieves the following two results by a gradual development: 1. The Ninth and the Tenth Commandments forbid selfishness, the root of every sin against neighborly love; 2. the Ninth and the Tenth Commandment forbid evil lust, the root of every sin. These results achieved, he uses them to establish the final theme: The Ninth and the Tenth Commandments constitute the summit of the whole Law. If, at the same time, he succeeded in influencing the emotions and the will, the effect was due, as in the case with many other catechists, to his personality rather than his method. He himself appears to have felt the insufficiency of his method; for he supplemented his “dialectic-didactic” method by another, which he
\quellseite{542}
called the “teleological-paraenetic” method. The latter was to be employed chiefly in the instruction of confirmands and in \emph{Christenlehre} when the teacher discusses Baptism, Confession, and Communion.

Just when the scientific catechization (\emph{Kunstkatechese}) had attained to dominating influence in catechetical instruction, a new idea began to gain ground. Sponsored by the psychology and pedagogy of Herbart, as developed especially by Ziller (p. 151), the thought gradually gained ascendancy that the way from the “concept” to the “conception” (cf. p. 205) is through the intuition, and that therefore a fundamental principle of a correct catechetical method must be first to feed the intuitive faculty of the child, and thus to enable him to make progress by independent thinking and judgment. Another factor making for a change in the method of instruction was the increasingly advocated fundamental idea of religious psychology that specifically religious conceptions can exert vitalizing power upon mind and soul only as there is a beholding—an intuition—of the religious and moral life, of which such conceptions are the expression. Thus the view at once was cast aside that linguistic usage or the Bible verse is to constitute the starting-point in the development of the truths of the Catechism. At the same time the scientific catechization (\emph{Kunstkatechese}) could no longer be advocated; for though intuitive material taken from history or the life of the present was employed in that method, the law governing it in all its parts was not the paramountcy of the intuitive faculty. The scientific catechization did not proceed from an appeal to the intuitive faculty; that is to say, it did not, first of all, strive to let the children behold the nature and meaning of the religious life upon the background of the actual life of religious-moral personalities, thus enabling them to recognize it in its worth, significance, and beauty, and thereby in their hearts to arouse interest in it, desire for it, and the determination to lead such a life themselves. And just because this was neglected
\quellseite{543}
by the advocates of the scientific catechization, their efforts were bound to be abstract, doctrinaire, anything but pulsating with life, and, in consequence, anything but vitalizing, unless the defects of the method were atoned for in large measure by the personality of the catechist. Whatever fault one may find with the school of Herbart and Ziller, and the modern religious philosophy, it must be clear from the whole second part of this book (especially pp. 223–226, but also 253, 260, etc.) that the principle: “from the intuition of the religious-moral life to the religious-moral conception” has to be a dominant and indispensable rule of the catechetical method. The catechist has to bring the children face to face with the characters of history, in order to help them to see what truth and Christian life are. This was the method of God the Father when He sent His Son into the world, so that men might behold in Him the Father, the Father’s nature and attributes, His holiness and love. While this was not the sole purpose of sending Christ into the world, it was nevertheless closely related to redemption as it was, one of the reasons of His coming (John 14:9; Matt. 11:27). This was the method of Christ, who said to John and Andrew: “Come and ye shall see!” (John 1:39, compare 46); who told the delegates of the Baptist: “Tell John the things which ye hear and see” (Matt. 11:4); who, in order to lead men to faith, did not formulate definitions or syllogisms concerning God, His own person, and the Kingdom of Heaven, but presented Himself in person, so that, beholding His divine-human life as a Redeemer and perceiving the powerful appeal of His personality, men might obtain a true vision, gain vital experiences, and form wholesome judgments (cf. John 4:42). This was the method of Paul also, who “painted” Christ before the eyes of the Galatians as the crucified (Gal. 3:1).

The catechetical rule here established has indeed been exaggerated, with a twofold error as result. There are
\quellseite{544}
those who think that all that is necessary now in religious instruction is to picture to the children the life experiences of grace-endowed personalities in fascinating concreteness and with “contagious” warmth, so that, coming under the spell of the life so beheld and mysteriously laid hold of by God in the depth of their own souls, they might “meet God”. It is evident that, where that view prevails, the emphasis is laid upon feeling, upon patterning one’s own feelings and experiences after the life of the personalities beheld; far less interest is evinced in the thoughts and truths uttered by them or shown forth in their lives; still less is it believed to be the object of religious instruction, in co-operation with the children and with the life of such grace-endowed personalities as object-lessons, to develop definite, precise, religious ideas and truths. In that case, the intuitive principle, by reason of gross exaggeration, has led to the complete abandonment of all instruction in the Catechism: the chief parts of the Christian faith would hardly be rated as anything higher than the expressions of the religious life of a Moses, an Athanasius, a Luther, and at best would be incidentally treated in connection with Biblical or Church History. That we cannot join anyone on such a path, is clear from ch. 26.—Others thought that the observance of this fundamental rule led of necessity to the abolition of instruction in the Catechism as an independent branch of religious instruction. The intuitive material for the religious-moral life, it was held, is found in the Bible stories. When, in the process of presentation, by a detailed description of the characters of the story an attractive, alluring picture has been painted of true religion and morality, which has been made even more precise, more clear, and more appealing in the process of penetration (and if need be, of comparison), nothing more remains to be done, it is concluded, than to gather the several features of the religious life that have been set forth, into
\quellseite{545}
a terse classical expression. This is supposed to be most readily accomplished by using for the purpose a hymn stanza, a passage of the Bible, or an appropriate sentence from the Catechism. Thus the conception, it was thought, grows out of the intuitive material. The catechetical thesis would not impress the children as something discovered and dragged in from the outside, to which they must submit at any price; but as something ascertained through their own effort in co-operation with the teacher, to which they cannot refuse to yield without self-disapproval. But, by the side of this “catechetical step” in instruction in Biblical History, it was held, there is neither room nor need nor justification for anything like specific instruction in the Catechism. In contrast to the view previously discussed there is here a recognition of the value to be accorded, also in the religious life, to clear cognitions and sharply defined truths and conceptions; and there is a lively interest in working these out. While this touches a sympathetic cord in one who is bound, with Paul and Luther, to lay stress upon “sound doctrine” (2 Tim. 4: 3), we know ourselves to be opposed also to this view (compare ch. 26). We believe that it is possible to hold to that fundamental rule in all its bearings and apply it in practise without the alternative of discarding specific instruction in the Catechism.

Two possibilities have to be considered in this connection: \emph{The first} is that the catechist treats Luther’s Catechism as something not yet existing; in every catechesis he proceeds from an intuitive basis; from this he gleans some religious truth or directs the children to glean it, garbs it in the words of the Catechism, and thus builds up the Catechism by a method purely analytical (this word to be understood in the sense of p. 442), until it stands as a unit before the eyes of the children. This process is likely to comprise the following steps in every catechetical lesson: 1. The
\quellseite{546}
Catechism truth is developed from some incident in the life of the child or in history; 2. When the children have received a clear idea of the subject matter, it is summarized and, with the aid of the teacher, first expressed in the child’s own words; 3. the teacher garbs such explanation in the form of a statement of Luther’s Catechism; 4. in this form it is written on the blackboard, drilled with due attention to proper expression; then—not before—it is looked up in the Catechism, drilled, and memorized. In this way, when, e. g., the Fifth Commandment is treated, with the story of Cain’s fratricide to serve as starting-point, the following truths can be established in an ascending line through the didactic conversation: 1. We are not permitted to hurt our neighbor in his body; 2. we are not allowed to harm our neighbor in his body; 3. we should help our neighbor in every bodily need; 4. we should befriend our neighbor in every bodily need; 5. all this we should do or leave undone out of fear and love of God. These sentences written on the blackboard, are then, at the close of the whole conversation, gathered into a whole in the words of the Catechism.—\emph{The second possibility} is that the Catechism is put in the children’s hands at the outset; the part to be considered is looked up; where necessary, its structure is set forth by a careful analytical process, in which the children take part; the analysis completed, its clauses are written on the blackboard, whereupon the children are made acquainted with the several sentences or clauses by means of the analytical development (as set forth on p. 442), which always rises from the intuition of religious-moral life to the religious-moral conception. If, for instance, the Fifth Commandment is the subject under consideration, let it be analyzed into its two component parts: 1. We should fear and love God and not hurt or harm our neighbor in his body; 2. we should fear and love God and help and befriend our neighbor in every bodily need. In
\quellseite{547}
connection with each one of the four truths contained in these two sentences, the catechist will make a beginning with the intuitive material offered by a certain Biblical story or the experience of the children, in order thus to lead them by the developing process from that point to a comprehensive understanding of the specific truth to be taught. A detailed picture of the slave labor of the Israelites in Egypt, for instance, may be painted, in order to let the children vividly see and feel the impossibility for anyone who fears and loves God to hurt or harm his neighbor in the body. Or a word-picture of the Good Samaritan, with due attention to detail, is painted, in order to let his noble traits arouse the sympathetic, moral, and religious interest of the pupils and to move them to the definite resolution to follow his example. Or if the Second Article is to be explained, Luther’s explanation is first analyzed, whereupon the several clauses may be developed in the way of induction (analytical), with an intuitive basis as starting-point in each case. In order to enable the child really to apprehend the truth “without Christ, man is a lost and condemned sinner”, a word-picture of Adam and Eve at the moment when God drove them from Paradise should be painted, or of some heathen from the history of Missions, whose conscience has been awakened, so that the children may read in the outlines of that picture, and in a measure reproduce in their own feelings, what the sentence from the Catechism expresses.

—There the Catechism, in its several clauses, and finally as a whole, is the result of the discourse; and here it is the starting-point; but in both cases the fundamental rule: from the intuition of the religious-moral life to the religious conception, dominates the process. In the first mode it lies at the root of the process as a whole; in the second it is of moment at least in the teaching of some specific truth of the Catechism. In both cases the important element of progress
\quellseite{548}
introduced by the method of rationalism as well as the scientific catechization is given due consideration; we have a real didactic conversation, in which the pupil is led from the unknown to the known, i. e., results are gained through diligent development.

Which of the two modes described deserves preference? The superiority of the first, at first glance, appears so great that there is no room for the second. All learning is based upon apperception (p. 212); for this reason all teaching of new truth, if it is to be correct methodically, is bound to take place in connection with such truth as is already found in the soul. This principle is fully observed in the first mode; for the starting-point is not some foreign quantity, as “Catechism”, but a personal experience or a well known story: therewith the cell-wall is supplied, as it were, against which the new cell is built. As a matter of fact, however, the case of the second mode is by no means so bad as might seem at first sight. It should be borne in mind that, for pupils who have been instructed in Biblical History according to the rules laid down above, the Catechism has ceased to be a “foreign quantity”. While, so far, it has not been presented to them as a unit or treated as such in special lessons, it is nevertheless true that every Biblical History period had taken up a certain part of the Catechism and garbed the truth gleaned from Biblical History in its language. It is not likely that there are many sentences of the Catechism which have not already been elaborated by the instruction in Biblical History and have not yet become familiar to the pupils. The need hardly exists at that juncture to construct the Catechism independently. On the contrary, the child is prepared for it, and desires to become acquainted with it as a whole. The child has by this time attained to such intellectual maturity that this life manual—intended and qualified to serve as fingerpost and compass
\quellseite{549}
for his whole religious and moral life, by the help of which he is to become more and more familiar with Holy Scripture—can safely be put into his hand as a summary of the faith confessed by the Church, full membership in which is the goal of his present tutelage. The second method is the very expedient calculated to impress upon the child the difference between the instruction in the Catechism now in progress and the instruction in Biblical History which preceded it and still continues as a parallel course. There is no doubt that the pupil will become more firmly grounded in the Catechism; for it is the Catechism which he now has always before him, and which now forms the starting-point for every catechesis and, properly understood, also its aim. Nor is it true that Christian children lose interest when the aim of the catechesis has been announced to them at the outset in the words of the Catechism; as, for instance, when a new general subject is announced in words such as these: “The next lessons shall teach us what Luther’s explanation of the First Commandment requires of us Christians”. The same holds good when some subordinate point is treated, as, “Today we want to learn how, as Christians, we should trust in God above all things.” There is the additional advantage here that the children are enabled at the outset to look upon the intuitive material to which they are introduced from the definite point of view indicated by the statement of aim, and that a close inner relation is at once established in their souls between the truth expressed in the Catechism and the intuitive material laid before them, so that the one cannot awake in the memory without the other. For this reason the second mode should receive preference as a rule, especially in the circumstances prevailing in our country. It adopts all the elements of truth contained in the catechetical methods heretofore employed and presents the correct method in a fusion of all of them. The
\quellseite{550}
catechist here presents the whole Catechism and holds it up to the young as the sum of truth, experienced and proved as such by himself and the whole mature congregation:—this is the element of strength in the method obtaining at the time of the Reformation. At the same time the catechist endeavors to analyze the several parts through questions, and thereby to bring about, or, in any event, facilitate, a rudimentary understanding of them:—this is the truth in the method of pietism. While constantly stimulating the pupil to independent thinking, he makes it his object to effect a real comprehension of the subject by means of the didactic conversation, and thus to put in operation the element of truth found in the scientific catechization as well as the Socratic method. But all this is done only by leading from the intuition of the religious-moral life to the religious-moral conception, whereby the compatibility of this method with the fundamental demands of present-day psychology stands demonstrated.

It remains to add one, final, requisite of correct method: The whole didactic conversation must pursue aims altogether practical. The needs to be satisfied should not merely be of an intellectual nature. While these are by no means to be overlooked, the molding of the life should invariably be the final object. Not this is of moment that the children, on the basis of the intuitive material supplied to them, are able to give a correct definition of “trusting”, but that, at the sight of David risking unequal combat with Goliath on the strength of his trust in God, they receive a vivid intuition of living faith; that their hearts are warmed by what they see; that they recognize the great value of trust in God; that they begin to perceive the utter inadequacy, the sinfulness of any other trust; that they recognize the fact that there are situations in their own life where everything depends upon just such trust; and that, if it please God
\quellseite{551}
(p. 238, 278), they leave the room with the resolution, even now as children, ever to trust in no one as they do in God. It is only thus that the personal character of Luther’s explanation with its “I”, “we”, “us”, comes into its own (p. 105); it is a direct challenge to the catechist not to content himself with the rôle of a mere preceptor, but by all means to be personal and practical in his instruction, and to keep in mind the individual and the molding of his life. Also here the threefold life of the soul is to be set in motion, which means that the intellect with the emotional and the volitional life must receive harmonious attention. Quite true, clear conceptions are to be achieved; for hazy ideas are seldom effective; but not these conceptions in themselves should be the aim, but only as a power for the practical Christian life. Luther may well serve as our pattern in this respect, who, no matter what important truths he states, e.g., in the First or Second Article, in conclusion always establishes the connection between them and the practical life.

Accordingly the process of a catechization (especially on the First, Second, and Third Chief Parts) will be as follows: \emph{(1) Statement of aim}, in order that the children collect their thoughts and concentrate them upon one point. The stage of preparation may usually be safely passed over, especially since the new catechesis is ushered in by review and drill questions concerning the substance of the preceding catechesis, whereby the children are led up to the point where a new start is made. The statement of aim should be couched in a single sentence of the Catechism, as, “Today let us learn that Christians should use the name of God for the purpose of praise and thanksgiving”; or, “Today let us learn that no one has ever become a believer in the Lord Jesus by his own reason and strength”. \emph{(2) The developing of the specific truth from the intuitive material}. If this truth
\quellseite{552}
is to be of value, result in clear cognitions, and move feelings and will, the intuitive material should not be merely touched upon: all its points should be utilized as far as they may contribute to the feature indicated in the statement of aim. Only thus can an adequate picture of God’s or man’s conduct be produced. There may be occasions when it will be necessary to bring in a second or a third story, in order to obtain all the features of the catechetical truth under consideration. When, for instance, trust in God, as required by the explanation of the First Commandment, is to be set forth, it will be advisable to adduce not only David’s combat with Goliath, but also his conduct under Saul’s persecution; for while the former story illustrates trust in action, the latter illustrates trust in suffering. The closer the relations between the story drawn upon, whether from the Bible or life, and the truth of the Catechism, the better: another reason, why nothing has been gained when a story is merely touched upon. \emph{(3) A comparison with similar or even contrasting stories} will often help to bring out more forcefully the truth under consideration. Thus David’s trust in God is rendered more conspicuous when compared with Goliath’s self-trust. Likewise, the impression made by David’s trust is deepened when paralleled with that of Moses which prompted him to pass through the Red Sea. God’s faithfulness appears so much greater when viewed against the background of man’s unfaithfulness. If the specific truth has been developed from the intuitive material and its true nature been more fully recognized by comparison, there follows, \emph{(4) a valuation} of the action of God or of such persons as have been used for intuitive material, or drawn upon for the purpose of comparison. The children are trained to rate the value and significance of an action in itself and in the light of its results; human acts are to be examined as to whether they are good or bad, worthy of emulation
\quellseite{553}
or of detestation. Thus has already been introduced what need not be emphasized as a special step in every case and must occasionally be omitted altogether, namely, \emph{(5) the application of the truth to the children’s own life}. Such application is effectual and in place only when it does not form a mere appendage but forces itself upon the mind as the natural conclusion of what has gone before. If, e.g., the high value of trust in God has been recognized, questions like the following will naturally suggest themselves: What shall you do when you are to solve a difficult problem; when you are to undertake a dangerous trip; when you are alone at night; when a severe thunderstorm is coming up; when you should lose your father; when poverty oppresses you; when you might become rich by leaving the right path, etc.? If the children, in connection with the Second Article, have been made conscious of the fact that Jesus Christ is true God, the question will naturally suggest itself: Who, then, will be your comfort in trouble, even in the greatest trouble; but what, in view of that fact, will you always have to guard against? If the truth has dawned upon the children that their Lord and Redeemer is true man also today—the same man who has gone through every temptation here on earth, they are shown how to apply to themselves the comfort given in the Epistle to the Hebrews, that the Saviour will surely have pity upon our weakness. If a description has been given of the price of our deliverance from sin, death, and devil, namely, Jesus’ blood and life, the question will at once occur how God evidently judges of our sin, and what care we should take to guard against it, etc. In connection with the fourth or the fifth step the Bible verse or the hymn may usually be drawn upon with advantage. In connection with the fourth, the Bible passage will usually, though not exclusively, come in as \emph{dictum probans}—in connection with the fifth usually as an admonition for life (p. 379).
\quellseite{554}
It is still indispensable as \emph{dictum probans}, modern objections to the contrary notwithstanding. For instance, when the redemptive value of the work of Christ is under discussion, who is authorized to render a final valuation in the premises save (in addition to the prophets) Jesus and the apostles? Who can be relied upon for a thorough and trustworthy disclosure of human sin and helplessness and of saving grace as well, man himself, ever ready to deceive himself and to trust in his own power for help, or God; the patient who believes himself much sounder than he is, or the physician, who understands the disease and is aware of the healing power of his remedies? To be sure, the Bible verse can be applied in a manner so mechanical that the children are led to feel that the Bible is merely a code of doctrine, to whose injunctions and statements they must blindly submit. But when, again and again, the truth is emphasized that, in a history of over two thousand years, the Bible has shown itself as truth and life; that it is the grace of God which has given it; that true association with it has never yet brought forth slaves but always free children of God—happy and cheerful in every trouble, joyous and bold even when face to face with death; when the children are made to feel that it is holy and precious to the catechist himself; that it is the true book of life, such thoughts will not arise. What the children need here, as in connection with instruction of children in general, is a little confidence in the greater experience of the teacher. Often it may be advisable already in connection with the fourth step, and, still more, with the fifth, to introduce passages from the Bible as testimonies expressing the experience of those holy men, provided only that there is no encroachment upon its absolute correctness and universal application. Cases in point are Ps. 23 in connection with the First Commandment and the First Article; Rom. 8:31 ff. when the blessed state of God’s children is under
\quellseite{555}
discussion; Ps. 145: 18. 19 in connection with the Second Commandment. Whatever Bible passages are called for under point five, should bear a practical character throughout, with Ps. 119: 9 and Ps. 19: 8—12 governing the perspective. When the time has come for the class to form a valuation of the act under discussion or to apply the truths derived, it may be appropriate to read lengthier sections from Scripture, as, Isa. 53 in connection with “purchased, redeemed, and won . . . with His holy precious blood, etc.”; Heb. 11 in connection with the treatment of the nature and the fundamental character of faith; 1 Cor. 13 in connection with the subject of Christian love of the brethren, etc. (p. 371). The hymn takes its place by the side of the Bible passage as the expression of the experience of the later generations and, at the same time, as evidence of the fact that also today men may have the same experiences which those ancient men of God experienced and expressed, under singular divine guidance, in Holy Scripture. There are occasions when whole hymns may be aptly quoted.

—The final requirement is \emph{(6) the drill}. If a catechism with explanation be the textbook, the section dealing with the subject is looked up when application is finished. It goes without saying that it will be of service only when it is the result or thesis of the catechization that had been conducted in the class-room. This section should then be read by the pupils at the close of the catechization. At home it should be studied. The next day, before the new subject is taken up, a number of questions are asked for the purpose or reviewing and drilling. The questions should not require answers that literally reproduce memorized material. It is better that merely the essential facts be recapitulated, as has been emphasized on p. 447, nor should they be restricted to intellectual features alone; see Practical Examples at the end of this section. If the catechism in use contains no
\quellseite{556}
explanation, a notebook should be procured by the pupils. In that the drill questions asked by the teacher should be entered and studied until the next lesson, when they are to be answered in writing or orally. For light upon this point, as well as for information how to give the catechization a practical trend, we refer to Practical Examples. At the same place is to be found a catechetical “excursus”, such as may be undertaken at the end of a Chief Part or after completing the Catechism in all its parts.

It has been emphasized above (p. 320) that the subject for catechization is not the Catechism text together with Luther’s explanation, but only the latter. In connection with the Ten Commandments it may possibly be necessary to enter briefly, and in a preliminary manner merely, upon the Old Testament significance in order subsequently to ascend to the lessons they contain for us Christians according to Luther. In connection with the Second Chief Part a verbal explanation of the text together with an outline will suffice; at most a few historic remarks are needed, so that afterward all available strength may be devoted to the setting forth of what we Lutheran Christians mean to confess in the Apostolic Creed, which is so often upon our lips. When, in the Third Chief Part, which requires but brief treatment, and in the Fourth and Fifth, which require a lengthier and somewhat different treatment from that of the first three, the biblical text is made the starting-point, the whole exposition will issue in Luther’s explanation.—When, finally, the requirement is made that the explanation of the Catechism should be \emph{Christo-centric}, all that is needed is to follow the trend of Luther’s explanation and the \emph{Christo-centric} character of the explanation will take care of itself. Cf. especially pp. 312 f., 324 ff., 351, but also 346.
